\section{Preparing the regular fibers}\label{sec:regular-preparation}

The torus fibration gives a way to average the pair along its regular
fibers.  We use this averaging to construct an invariant pair and a
closed form of unit fiber area orthogonal to both members.  These are
the data needed for the auxiliary volume equation in
\Cref{sec:auxiliary-volume}.  Every change to the pair is exact, and the
curve criterion keeps a taming form in the original third class
throughout the preparation.

Write $p\colon X\to S$ for the fibration of \Cref{thm:pencil}, and let
$\Delta\subset S$ be its finite set of critical values.  Put
\[
 S^\circ=S\setminus\Delta,
 \qquad X^\circ=p^{-1}(S^\circ).
\]
The fiber and base orientations are the complex orientations supplied by
that theorem.  We identify the class of an oriented fiber with its
Poincar\'e dual $F\in H^2(X;\Z)/\mathrm{torsion}$.  In particular, if
$\nu_S$ is a positive area form with $\int_S\nu_S=1$, then
\begin{equation}\label{eq:regular-basic-class}
 [p^*\nu_S]=F.
\end{equation}
For example, this follows by first taking a representative of the unit
class supported near a regular value, whose pullback is a Thom form for
the corresponding fiber.

\begin{proposition}\label{prop:regular-preparation}
Let $(A_0,C_0,B_0)$ and $p\colon X\to S$ be the positive triple and
fibration obtained in \Cref{thm:pencil}.  Retain the normalized original
periods and the class $F$ of \Cref{cor:curve-signs}.  On disjoint disks
$U_j$ about the critical values, use the models of
\Cref{thm:nodal-model}, with
\begin{equation}\label{eq:regular-input-model}
 \Sigma_j=h_j\,\dd s\wedge\dd z,
 \qquad A_0=\Re\Sigma_j,
 \qquad C_0=x_j\Re\Sigma_j+y_j\Im\Sigma_j,
 \qquad y_j\ne0,
\end{equation}
where $h_j\ne0$, $x_j$, and $y_j$ are constants and
$z=a+\tau(s)b$ modulo the elliptic period lattice.

There are smaller disks $U'_j\Subset U_j$ and a smooth path of positive
closed triples $(A_t,C_t,B_t)$, $0\le t\le1$, starting at
$(A_0,C_0,B_0)$, with the following properties.
\begin{enumerate}[label=\textup{(\roman*)}]
\item Each of the three cohomology classes is constant.  The restrictions
of $A_t,C_t$ to the fibers vanish, and the pair equals $(A_0,C_0)$ on
$p^{-1}(U'_j)$ for every $t$.
\item The endpoint pair $A=A_1$, $C=C_1$ is invariant under the torus
group bundle acting on $X^\circ$ by the action-angle translations of
$A_0$.  This is also the torus action determined by $A$.
\item There is a smooth closed invariant two-form $\Psi$ on $X^\circ$
such that
\begin{equation}\label{eq:regular-prepared-data}
 \int_{p^{-1}(s)}\Psi=1,
 \qquad \Psi\wedge A=\Psi\wedge C=0
 \quad\text{on }X^\circ.
\end{equation}
On each punctured $p^{-1}(U'_j\setminus\{s_j\})$ it equals
$\dd a\wedge\dd b$.  In particular, $\Psi^2$ vanishes there.
\end{enumerate}
The fibration and its smaller nodal models remain fixed.  All of these
choices are made before introducing the parameter $\eps$ in the
auxiliary-volume construction.
\end{proposition}

We first describe averaging and the positive pair paths it permits.

\subsection{Fiber translations and exact averaging}

The regular torus fibration is a torsor: there are local choices of an
origin in each fiber, but a global choice is unnecessary.  Here is a
direct description of its translations.  If $\alpha_s\in T_s^*S^\circ$,
the equation
\[
 \iota_{Y_\alpha}A_0=-p^*\alpha
\]
defines a unique vertical vector field along $p^{-1}(s)$, since the fibers
are Lagrangian for $A_0$.  Locally take $\alpha=\dd f$ on the base.  The
fields $Y_{\dd f}$ are Hamiltonian; their brackets vanish because their
Hamiltonians are pulled back from the base and their fields are vertical.
Their complete flows on each compact fiber therefore give a transitive
$T_s^*S^\circ$-action with a rank-two stabilizer lattice $\mathcal L_s$.
The lattices vary smoothly locally: choose a local origin section and
apply the implicit function theorem to the return equation for each
member of a lattice basis.  The derivative in the orbit parameter is
an isomorphism onto the vertical tangent plane, so the two periods
continue smoothly and remain a lattice basis after shrinking the base
chart.

A local smooth section $\alpha$ of $\mathcal L$ is a closed one-form.
Indeed, its vertical flow $\phi_t$ has $\phi_1=\Id$, whereas Cartan's
formula gives
\[
 \phi_t^*A_0=A_0-tp^*(\dd\alpha).
\]
Putting $t=1$ proves $\dd\alpha=0$.  A local lattice basis thus consists
of closed one-forms, and these are differentials of local base
coordinates.  The corresponding period-one flows define translations
by constant angles in $\R^2/\Z^2$, and each such translation preserves
$A_0$.  In two overlapping choices of angle coordinates the change has
the form
\begin{equation}\label{eq:regular-angle-transition}
 \theta'=M\theta+g(s),\qquad M\in\operatorname{GL}(2,\Z)
 \text{ locally constant}.
\end{equation}
Thus translation by $u$ in one chart is translation by $Mu$ in the
other, independently of the base-dependent change of origin $g$.

On a model disk, the translations are exactly
$z\mapsto z+a_0+\tau(s)b_0$ for constant real $a_0,b_0$ modulo integers.
They preserve $\Sigma_j$, since the extra term in $\dd z$ is a multiple
of $\dd s$.  Equivalently, contraction of $A_0$ with the period-one
fields $\partial_a,\partial_b$ gives the closed base one-forms
$-\Re(h_j\dd s)$ and $-\Re(h_j\tau(s)\dd s)$.  They form the period
basis just described.  Both model forms in
\Cref{eq:regular-input-model} are consequently invariant.

\begin{lemma}\label{lem:regular-averaging}
There is a well-defined averaging projection
$\operatorname{Av}\colon\Omega^*(X^\circ)\to\Omega^*(X^\circ)$ onto
forms invariant under these local translations, and it commutes with
$\dd$.  If a closed form $\xi$ satisfies $\operatorname{Av}\xi=0$, then
$\xi=\dd\beta$ for a smooth form $\beta$ with
$\operatorname{Av}\beta=0$.  The primitive can be chosen without
enlarging support over the base: if $\xi$ vanishes over an open set in
$S^\circ$, then so does $\beta$.

In particular, a closed form and its average represent the same class
on $X^\circ$.  If their difference is supported over a compact subset
of $S^\circ$, the primitive extends by zero over the nodal neighborhoods
and gives exactness on $X$.
\end{lemma}

\begin{proof}
In a local translation chart define the average by integrating pullbacks
against normalized Haar measure on $\R^2/\Z^2$.  The conjugacy in
\Cref{eq:regular-angle-transition} and invariance of Haar measure make
the definitions agree on overlaps.  Each pullback commutes with $\dd$,
which proves the assertion about differentiation.

For exactness, choose a translation-invariant horizontal splitting of
$TX^\circ\to p^*TS^\circ$ and a smoothly varying invariant flat metric
on the vertical tori.  Such choices exist: make them locally and patch
horizontal lifts and vertical inner products by partitions of unity
pulled back from the base.  The transition rule above preserves
invariance, and the relevant sets of choices are affine or convex.

Filter differential forms by horizontal degree.  On the associated
graded complex the leading differential is the vertical exterior
derivative $\dd_v$.  For each horizontal degree, fiberwise Hodge theory
\cite[Equations~(1.11)--(1.22)]{TaylorHodge2010}
provides a homotopy $H_v=\dd_v^*G_v$ on the vertical complex, with the
usual sign for its horizontal degree.  Here $G_v$ is the inverse of the
vertical Laplacian on the complement of its kernel.  Its harmonic forms
are precisely the translation-invariant forms.  On the zero-average
subspace, therefore,
\begin{equation}\label{eq:regular-vertical-homotopy}
 \dd_v H_v+H_v\dd_v=\Id.
\end{equation}
These operators act smoothly in the base: in a local torus chart they
are the inverses on nonzero Fourier modes for a smoothly varying flat
metric.  They preserve zero average and use no values at other base
points.

Suppose the smallest horizontal degree of the closed zero-average form
$\xi$ is $r$, and denote its component of that degree by $\xi_r$.
Closedness gives $\dd_v\xi_r=0$.  By
\Cref{eq:regular-vertical-homotopy}, subtracting
$\dd(H_v\xi_r)$ removes $\xi_r$, and the remainder is closed, has zero
average, and has horizontal degree at least $r+1$.  Repeat this step.
The horizontal degree is at most two, so the procedure terminates.  A
component of vertical degree zero that is vertically closed and has
zero average is already zero.  The sum of the primitives gives $\beta$.
Every operation is local in the base, including the differentiations
used in the successive remainders.  Hence the support assertion follows.
Applying the construction to $\xi-\operatorname{Av}\xi$ proves the
remaining claims.
\end{proof}

\subsection{The positive cone and the third class}

\begin{lemma}\label{lem:regular-mixed-cone}
Fix a symplectic two-form $A$ with $A^2>0$ and a Lagrangian two-plane $L$ in an
oriented four-dimensional vector space.  Among two-forms $C$ satisfying
$C|_L=0$, each component of the condition that $(A,C)$ be a definite
pair is convex.  Adding a two-form pulled back from the quotient by
$L$ to either member of the pair does not change its wedge-product
matrix.

Consequently the path from $C$ to its translation average, with $A$
fixed, preserves definiteness for a fiber-Lagrangian pair on $X^\circ$.
Basic additions to either member do so as well.  Along these paths the
fiber complex orientation and the quotient complex orientation remain
the ones fixed at the start.
\end{lemma}

\begin{proof}
Choose covectors $e_1,e_2$ vanishing on $L$ and complementary covectors
$f_1,f_2$ so that
\[
 A=e_1\wedge f_1+e_2\wedge f_2,
 \qquad
 C=\sum_{i,j=1}^2 M_{ij}e_i\wedge f_j+c\,e_1\wedge e_2.
\]
Direct wedge multiplication yields
\[
 A\wedge C=\tfrac12\tr(M)A^2,
 \qquad C^2=\det(M)A^2.
\]
Writing $h=\tr(M)/2$ and $T=M-h\Id$ gives
\begin{equation}\label{eq:regular-mixed-square}
 (rA+tC)^2=\bigl((r+th)^2+t^2\det T\bigr)A^2.
\end{equation}
Thus the pair is definite exactly when $\det T>0$.  Set
\[
 T=\begin{pmatrix}u&v\\w&-u\end{pmatrix},\qquad
 q=\tfrac12(v-w),\qquad r_0=\tfrac12(v+w).
\]
Then $\det T=q^2-r_0^2-u^2$.  Its two positive components are
$q>\sqrt{r_0^2+u^2}$ and $q<-\sqrt{r_0^2+u^2}$, both convex; $h$ and
$c$ are unrestricted.  This proves the first assertion.  A quotient
two-form wedges to zero with both members of a fiber-Lagrangian pair,
and its own square vanishes.  This proves the assertion about basic
additions, also when both members change.

At a fixed point all translated pullbacks of $C$ lie in one of the two
components: translations preserve $A$ and can be connected to the
identity within the torus.  Their compact average remains strictly
inside that convex component, as does the segment to the average.
The common Lagrangian plane is a complex line by
\Cref{lem:definite-pair}.  Continuing the choice of sign of the almost
complex structure along the path keeps its orientation on this plane,
and hence the quotient orientation, fixed.
\end{proof}

The next lemma separates the smooth choice of representatives from the
pointwise question of whether a class contains a tamer.  It will also
be used when the final endpoint representatives have been chosen.

\begin{lemma}\label{lem:smooth-tamers}
Let $J_t$, $0\le t\le1$, be a smooth family of almost-complex
structures on a compact manifold. Suppose each $J_t$ admits a closed
tamer in one fixed class $H$. Given such representatives $\beta_0$
and $\beta_1$ at the endpoints, there is a smooth family of closed
tamers $\beta_t\in H$ with those endpoints.
\end{lemma}

\begin{proof}
A fixed tamer for $J_{t_0}$ remains a tamer on a parameter neighborhood
of $t_0$, by openness and compactness of $X$. Choose finitely many
such neighborhoods and corresponding fixed representatives. Near
$0$ and $1$ use the prescribed representatives. Choose a subordinate
smooth nonnegative partition of unity in the parameter, with the
endpoint representative having weight one sufficiently near its
endpoint. Their weighted sum is closed on $X$ and represents $H$,
because all coefficients depend only on $t$ and sum to one. It tames
$J_t$ by convexity of the taming cone.
\end{proof}

\begin{lemma}\label{lem:regular-tamer-path}
Suppose $(A_t,C_t)$ is a smooth compact path of definite closed pairs
with the original two periods.  Suppose the fibers of $p$ remain common
Lagrangians with the original complex orientation, and the pair equals $(A_0,C_0)$ on smaller nodal neighborhoods.  Then the class $[B_0]$
contains a taming form for the associated almost complex structure at
every parameter.  If the initial pair is $(A_0,C_0)$, these representatives
can be chosen smoothly starting at $B_0$; an endpoint representative
can also be prescribed whenever it tames the endpoint structure.
\end{lemma}

\begin{proof}
Let $J_t$ be the structures defined by the pair, with the continued
sign.  First produce a convenient taming class.  On the regular region,
$J_t$ preserves the vertical tangent planes and induces an oriented
complex structure on each quotient $T_xX/\ker\dd p_x$.  This quotient
structure need not be independent of the point in the fiber.  Nevertheless
the positive base form satisfies
\[
 p^*\nu_S(\zeta,J_t\zeta)>0
 \quad\text{whenever }\dd p(\zeta)\ne0,
\]
and it vanishes on vertical vectors.  The restriction of $B_0$ to the
oriented vertical planes stays positive.

On a compact regular region containing the supports of the pair
changes, use a smooth splitting $\zeta=v+h$ into vertical and
horizontal vectors and any fixed background norm.  Compactness in
space and parameter gives positive constants $a,b$ and a finite $M$
such that
\begin{align*}
 B_0(\zeta,J_t\zeta)&\ge a|v|^2-M|v||h|-M|h|^2
                       \ge \tfrac a2|v|^2-M'|h|^2,\\
 p^*\nu_S(\zeta,J_t\zeta)&\ge b|h|^2.
\end{align*}
Thus $B_0+Kp^*\nu_S$ tames for one sufficiently large $K>0$.
On the unchanged nodal neighborhoods, $B_0$ already tames and
$p^*\nu_S$ is semipositive.  Indeed, the model matching in
\Cref{lem:pencil-insertion} chooses its complex structure with the
sign tamed by $B_0$, and its projection is holomorphic for that sign.
The same $K$ therefore works on all of $X$ and at every parameter.

Put $H_K=[B_0]+KF$ and
$V=[A_0]^\perp\cap[C_0]^\perp$.  By
\Cref{eq:regular-basic-class}, $H_K$ is the class just constructed.
It belongs to $V$, and
\[
 H_K^2=1+2K[B_0]F>0,\qquad
 H_K[B_0]=1+K[B_0]F>0.
\]
Consequently $H_K$ and $[B_0]$ lie in the same timelike component of
$V$.  To apply \Cref{thm:taming-cone}, let $d$ be the class, identified
with its Poincar\'e dual, of a nonconstant irreducible $J_t$-curve.
It has $H_Kd>0$.  By \Cref{cor:curve-signs}, either $Fd\ne0$ and its
sign is the sign of the coefficient $c$ in $d^+=cF^+$, or $d=cF$.
Also $[B_0]d=c[B_0]F$, with $[B_0]F>0$.  If $[B_0]d$ were
nonpositive, the alternatives would give $c<0$ and $Fd\le0$, so
$H_Kd=[B_0]d+KFd<0$, a contradiction.  Thus every such curve pairs
positively with $[B_0]$.  The hypotheses of
\Cref{thm:taming-cone} hold with the already tamed class $H_K$, and that
theorem supplies a tamer in $[B_0]$.

When the initial pair is $(A_0,C_0)$, apply \Cref{lem:smooth-tamers}
to the family $J_t$ and the fixed class $[B_0]$, using $B_0$ at the
initial endpoint and any prescribed endpoint tamer.  This gives the
asserted smooth representatives.
By \Cref{lem:definite-pair}, their triples with $(A_t,C_t)$ are positive.
\end{proof}

\subsection{Construction of the prepared pair}

We have established that averaging and basic additions preserve the
definite pair, and that the original third class remains available
along the resulting paths.  It remains to arrange the two orthogonality
equations for the unit-fiber-area form without changing the nodal data.

\begin{proof}[Proof of \Cref{prop:regular-preparation}]
We first average the pair and construct a closed invariant form of
unit fiber area.  We then adjust its two total pairings so that exact
basic changes of the pair can impose pointwise orthogonality.
Throughout, disks may be replaced by smaller nested disks, with
closures inside the original model disks.

\emph{Averaging the second form.}
On $X^\circ$ set $\overline C=\operatorname{Av}C_0$.  It is closed
and equals $C_0$ on every model disk, so it extends smoothly over $X$.
The difference $\overline C-C_0$ is supported over a compact subset
of $S^\circ$ and is exact on $X$ by
\Cref{lem:regular-averaging}.  The path
\[
 (A_0,(1-t)C_0+t\overline C)
\]
is a definite pair path by \Cref{lem:regular-mixed-cone}, with the same
periods and the same oriented fiber planes.  It admits a completion in
$[B_0]$ by \Cref{lem:regular-tamer-path}.  For the remaining construction
write $A=A_0$ and $C=\overline C$.  Both are invariant.

\emph{A closed form of unit fiber area.}
The positive number $q_0=[B_0]F$ is the common integral of $B_0$ over
the regular fibers.  Start with
\[
 \Psi^{(0)}=\operatorname{Av}(B_0/q_0)
 \quad\text{on }X^\circ.
\]
This is closed and invariant, and its restriction to every fiber is
the invariant two-form of integral one.

We next make it equal to $\dd a\wedge\dd b$ near every puncture.
On a punctured model disk, the two real two-homology generators are
the fiber and the torus obtained by sweeping the invariant $a$-cycle
around a base circle at $b=0$.  One can see this from the homology
sequence for a torus mapping torus: the monodromy is one shear, so
the fiber contributes one generator and the invariant part of its
first homology contributes one more.  The second generator bounds a
three-chain in the full nodal neighborhood, as described in
\Cref{thm:nodal-model}.  The smooth closed form $B_0/q_0$ consequently
has periods $1$ and $0$ on these generators.  Averaging does not change
them, by \Cref{lem:regular-averaging}.  The form
$\dd a\wedge\dd b$ has the same periods: its fiber integral is one,
and its restriction to the sweep at $b=0$ is zero.  It is a well-defined
closed form under the shear transition of the angular coordinates.

It follows that
$\Psi^{(0)}-\dd a\wedge\dd b=\dd\beta_j$ on the punctured
neighborhood.  Replace $\beta_j$ by its average, so it is invariant.
Choose a basic cutoff $\chi_j$ that is one on a smaller punctured
disk and zero near the outer boundary, and replace $\Psi^{(0)}$ there
by
\[
 \Psi^{(0)}-\dd\bigl((p^*\chi_j)\beta_j\bigr).
\]
These disjoint modifications give a smooth closed invariant form
$\Psi^{(1)}$ on $X^\circ$, equal to $\dd a\wedge\dd b$ near the
punctures.  Its fiber restriction is unchanged: the two invariant
forms being compared already have the same restriction of integral
one, and a differential of a basic cutoff has zero vertical
restriction.

\emph{The obstruction to an exact orthogonality correction.}
The form $\Psi^{(1)}$ has the required fiber restriction and nodal
behavior.  It remains to arrange orthogonality to the pair while
preserving the pair's cohomology classes.  Invariance reduces these
pointwise equations to a calculation on the base.

Let $\Phi$ be an invariant two-form on $X^\circ$ whose restriction
to each fiber has integral one.  For any invariant four-form $\Xi$,
we have the pointwise identity
\begin{equation}\label{eq:regular-invariant-top-form}
 \Phi\wedge p^*(p_*\Xi)=\Xi.
\end{equation}
Indeed, only the vertical restriction of $\Phi$ contributes to its
wedge with a basic two-form.  Invariant top forms have constant
vertical density on each fiber and are therefore determined by their
fiber integral; both sides of
\Cref{eq:regular-invariant-top-form} have the same integral on every
fiber after fixing two base tangent vectors.

In particular, replacing $A$ by
$A-p^*p_*(\Phi\wedge A)$ makes it orthogonal to $\Phi$, and the
same holds for $C$.  On the closed oriented surface $S$, a smooth
two-form is exact precisely when its integral is zero.  We therefore
first adjust $\Psi^{(1)}$ so that both resulting base two-forms have
zero integral.  Their vanishing near the punctures will allow us to
regard them as smooth forms on all of $S$.

\emph{Removing the two total pairings.}
There are smooth closed invariant forms $\widetilde A,\widetilde C$
on $X$, representing $[A],[C]$, supported over a compact subset of
$S^\circ$, and vanishing on vertical planes.  To construct them,
observe that a full one-node neighborhood has the homotopy type of a
torus with its primitive vanishing cycle collapsed; its real second
homology is generated by the fiber.  Hence $A,C$ are exact there,
since their fiber integrals vanish.  Subtract differentials of local
primitives multiplied by basic cutoffs equal to one near the nodal
fiber.  This gives global cohomologous forms that vanish on smaller
model neighborhoods.  Their restrictions to vertical planes remain
zero, because both the original forms and the differentials of the
cutoffs have that property in the required restriction.  Averaging
these forms, and extending by zero near the nodes, gives
$\widetilde A,\widetilde C$.  Their cohomology classes are unchanged
by \Cref{lem:regular-averaging}.

On a punctured model disk, $\Sigma_j=h_j\dd s\wedge(\dd a+\tau\dd b)$:
the term $b\tau'(s)\dd s$ in $\dd z$ disappears after wedging with
$\dd s$.  Thus $A,C$ have only mixed base--fiber terms there, and
\[
 (\dd a\wedge\dd b)\wedge A
 =(\dd a\wedge\dd b)\wedge C=0.
\]
Consequently the integrals
\[
 I_A=\int_{X^\circ}\Psi^{(1)}\wedge A,
 \qquad
 I_C=\int_{X^\circ}\Psi^{(1)}\wedge C
\]
are well-defined: their integrands vanish over smaller punctured
disks.  The matrix of pairings of $\widetilde A,\widetilde C$ with
$A,C$ is their cohomology Gram matrix, hence the identity under our
normalization.  Set
\begin{equation}\label{eq:regular-psi-class-correction}
 \Psi=\Psi^{(1)}-I_A\widetilde A-I_C\widetilde C.
\end{equation}
The form $\Psi$ remains closed, invariant, of unit fiber area, and
equal to $\dd a\wedge\dd b$ near the punctures.  By construction,
\begin{equation}\label{eq:regular-zero-total-pairings}
 \int_{X^\circ}\Psi\wedge A
 =\int_{X^\circ}\Psi\wedge C=0.
\end{equation}
Without the normalization one would solve the same two equations
using the positive invertible period Gram matrix.

\emph{Removing the pointwise pairings.}
Fiber integration of the invariant four-forms gives base two-forms
\[
 \zeta_A=p_*(\Psi\wedge A),\qquad
 \zeta_C=p_*(\Psi\wedge C).
\]
They vanish near the punctures and so extend by zero to smooth forms
on $S$.  Their integrals are zero by
\Cref{eq:regular-zero-total-pairings}.  Since $S$ is a closed connected
oriented surface, there are smooth one-forms $\kappa_A,\kappa_C$ on
$S$ with
$\zeta_A=\dd\kappa_A$ and $\zeta_C=\dd\kappa_C$.

Now make the simultaneous basic deformation
\begin{equation}\label{eq:regular-final-basic-path}
 A_t'=A-tp^*\zeta_A,
 \qquad C_t'=C-tp^*\zeta_C,
 \qquad 0\le t\le1.
\end{equation}
Both changes are globally exact, since the primitives are the smooth
pullbacks of $\kappa_A,\kappa_C$.  The changed forms are invariant,
have zero vertical restriction, and agree with the original pair
on smaller model neighborhoods.  Their wedge-product matrix is
unchanged by \Cref{lem:regular-mixed-cone}, so the path remains
definite.  \Cref{eq:regular-invariant-top-form} shows that its endpoint
satisfies
\[
 \Psi\wedge A_1'=\Psi\wedge C_1'=0.
\]
Moreover, contraction of a basic two-form with a vertical vector is
zero.  Hence the vertical Hamiltonian fields of base functions for
$A_t'$ are the same as those for $A$, and their period lattice and
torus translations are unchanged.

Finally concatenate the two pair paths, making their parameters flat
at the join.  They preserve the original pair periods, oriented
fibers, and smaller nodal models.  \Cref{lem:regular-tamer-path} gives
a smooth completion by forms in $[B_0]$, with initial value $B_0$.
This proves all the assertions.  We henceforth denote the endpoint
pair again by $A,C$ and its chosen third form by $B$.
\end{proof}
